% ============================================================================
% Brazilian Portuguese UI and structural strings for the One Chemistry Book series.
% Loaded by styles/onechemistry.sty when \booklang is pt.
% ============================================================================

% Theorem / statement environment titles
\newcommand{\omnameDefinition}{Definição}
\newcommand{\omnameTheorem}{Teorema}
\newcommand{\omnameProposition}{Proposição}
\newcommand{\omnameLemma}{Lema}
\newcommand{\omnameCorollary}{Corolário}
\newcommand{\omnameMethod}{Método}
\newcommand{\omnameExample}{Exemplo}
\newcommand{\omnameNotation}{Notação}
\newcommand{\omnameRemark}{Observação}
\newcommand{\omnameExercise}{Exercício}
\newcommand{\omnameExercises}{Exercícios}
\newcommand{\omnameProblem}{Problema}
\newcommand{\omnameProblems}{Problemas}
\newcommand{\omnameProof}{Demonstração}

% Admitted results and solutions appendix headers
\newcommand{\omadmittedtext}{Admitido neste nível.}
\newcommand{\omsolutionof}{Solução de}
\newcommand{\ompage}{página}
\newcommand{\omnameGotoSolution}{Solução p.}

% Misc text macros
\newcommand{\st}{\text{ tal que }}

% Cover / title page
\newcommand{\bookmaintitle}{One Chemistry Book}
\newcommand{\booktagline}{O livro de química para governá-los a todos}
\newcommand{\bookdescription}{A química do primeiro ano escolar ao fim da graduação}
\newcommand{\bookcontributors}{Os colaboradores do One Chemistry Book}
\newcommand{\bookversionlabel}{Versão}

% Colophon
\newcommand{\omcolophonsource}{Código-fonte, issues e contribuições:}
\newcommand{\omcolophonlicense}{Salvo indicação em contrário, este livro é distribuído sob a licença Creative Commons CC BY-NC-SA 4.0:}
\newcommand{\omcolophoncredit}{Cite como: One Chemistry Book, One Course (one-course.com).}

% Chemistry boxes (styles/onechemistry.sty): the recall box that opens a
% chapter building on another one, the lab box, the history box, the safety
% box. Placeholder wording by the coordinator (2026-10-04); the Book 1
% `pt` agent owns it and settles the register (the recall title addresses
% the reader in the same register as the book's exercise stems).
\newcommand{\omnameRecall}{O que você já sabe}
\newcommand{\omnameInTheLab}{No laboratório}
\newcommand{\omnameHistory}{História}
\newcommand{\omnameSafety}{Segurança}

% Solutions appendix part title
\newcommand{\omsolutionspart}{Soluções dos exercícios}

% Part titles (Anos 1–12)
\@namedef{omstr@part.grade1}{Ano 1 (6--7 anos)}
\@namedef{omstr@part.grade2}{Ano 2 (7--8 anos)}
\@namedef{omstr@part.grade3}{Ano 3 (8--9 anos)}
\@namedef{omstr@part.grade4}{Ano 4 (9--10 anos)}
\@namedef{omstr@part.grade5}{Ano 5 (10--11 anos)}
\@namedef{omstr@part.grade6}{Ano 6 (11--12 anos)}
\@namedef{omstr@part.grade7}{Ano 7 (12--13 anos)}
\@namedef{omstr@part.grade8}{Ano 8 (13--14 anos)}
\@namedef{omstr@part.grade9}{Ano 9 (14--15 anos)}
\@namedef{omstr@part.grade10}{Ano 10 (15--16 anos)}
\@namedef{omstr@part.grade11}{Ano 11 (16--17 anos)}
\@namedef{omstr@part.grade12}{Ano 12 (17--18 anos)}

% Part titles (Graduação 1–3)
\@namedef{omstr@part.bachelor1}{Graduação 1\texorpdfstring{\textsuperscript{o}}{o} ano (18--19 anos)}
\@namedef{omstr@part.bachelor2}{Graduação 2\texorpdfstring{\textsuperscript{o}}{o} ano (19--20 anos)}
\@namedef{omstr@part.bachelor3}{Graduação 3\texorpdfstring{\textsuperscript{o}}{o} ano (20--21 anos)}

\newcommand{\omnameDefinitions}{Definições}
\newcommand{\omnameTheorems}{Teoremas}
\newcommand{\omnamePropositions}{Proposições}
\newcommand{\omnameLemmas}{Lemas}
\newcommand{\omnameCorollarys}{Corolários}
\newcommand{\omnameMethods}{Métodos}
\newcommand{\omnameExamples}{Exemplos}
\newcommand{\omnameNotations}{Notações}
\newcommand{\omnameRemarks}{Observações}
\newcommand{\omnameChapter}{Capítulo}
\newcommand{\omnameChapters}{Capítulos}
% Section, for \cref{sec:…} (2026-10-04: it had no \crefname, so cleveref
% printed its English built-in "Section" in every edition).
\newcommand{\omnameSection}{Seção}
\newcommand{\omnameSections}{Seções}

% siunitx and cleveref keep their English conjunctions unless told otherwise:
% \qtyrange printed "5 to 4186 Hz" and \cref{a,b} printed "Capítulos 25 and 29"
% in the middle of Portuguese prose.
% cleveref installs its conjunctions at \begin{document}, so redefine them there.
% siunitx typesets range phrases and list separators in MATH MODE, where a
% bare string loses its spaces (\qtyrange inside $...$ printed "1a100") and a
% non-ASCII character is worse: "a" expands to \`a, which LaTeX reports as
% "Command \` invalid in math mode" and drops from the page. siunitx's own
% defaults are wrapped in \text{} for exactly this reason; these overrides
% must be too. Found on Biology Book 3 fr, 2026-09-05, where the book's 56
% \qtyrange sites made it visible; it was latent in every language edition of
% every book in this repo.
\sisetup{range-phrase={\text{ a }}, list-pair-separator={\text{ e }},
         list-final-separator={\text{ e }}}
\AtBeginDocument{%
  \renewcommand{\crefpairconjunction}{ e }%
  \renewcommand{\creflastconjunction}{ e }%
  \renewcommand{\crefrangeconjunction}{ a }%
  % a \cref spanning two label types builds a group list, which has its
  % own pair of conjunctions: "Teoremas 2.14 e 2.18 and Lema 2.17".
  \renewcommand{\crefpairgroupconjunction}{ e }%
  \renewcommand{\creflastgroupconjunction}{ e }%
}

% \today is English without babel (which is optional here), so the cover date
% read "July 27, 2026" on a Portuguese book. Brazilian Portuguese writes it as
% "27 de julho de 2026", month lowercase.
\renewcommand{\today}{\number\day~de~\ifcase\month\or janeiro\or fevereiro\or
  março\or abril\or maio\or junho\or julho\or agosto\or setembro\or outubro\or
  novembro\or dezembro\fi\space de~\number\year}

\renewcommand{\chaptername}{Capítulo}
\renewcommand{\partname}{Parte}
\renewcommand{\contentsname}{Sumário}
\renewcommand{\appendixname}{Apêndice}
\renewcommand{\indexname}{Índice}

% Periodic-table legend (\omperiodictable[legend], styles/onechemistry.sty).
% Coordinator wording, 2026-10-04 (it used to be hard-coded English in the
% style file, so no edition could translate it); the Book 1 agent of this
% language checks it against the book's own terminology.
\newcommand{\omnamePTmetal}{metais}
\newcommand{\omnamePTmetalloid}{metaloides}
\newcommand{\omnamePTnonmetal}{não metais}
\newcommand{\omnamePTs}{bloco s}
\newcommand{\omnamePTp}{bloco p}
\newcommand{\omnamePTd}{bloco d}
\newcommand{\omnamePTf}{bloco f}
\newcommand{\omnamePTalkali}{metais alcalinos}
\newcommand{\omnamePTalkaline}{metais alcalinoterrosos}
\newcommand{\omnamePThalogen}{halogênios}
\newcommand{\omnamePTnoble}{gases nobres}
